\section{Further research questions and a proposed programme}
\label{sec:agenda}

The strongest next steps are those that convert the new structural
information into a sharper theorem or a finite exact identity
certificate. The following questions are separated by their mathematical
obstruction and the kind of evidence that would count as progress.

\subsection{Sharp fixed-pair Euler errors}

For fixed $a,b>0$, determine
\[
 C(a,b)=\max_{N\ge1}R_N(a,b)
\]
and the set of maximizing indices. Existence follows from positivity and
$R_N\to0$, but the maximizer can occur after the first remainder.
Does the sequence have at most one local maximum outside the parameter
regions where monotonicity is already proved? Are there curves in the
$(a,b)$ plane across which the maximizing integer changes? A useful
first goal is an analytic classification of the sign of $R_{N+1}-R_N$
using the two-variable positive probability kernel.

The limit $a\downarrow0$ at fixed $b$ gives $C(a,b)\to C(b)$ by the
first-index lower bound and the uniform upper bound. Determine
a uniform first correction in $a$, and whether the maximizing index is
eventually $1$ on compact $b$ intervals. Joint limits where $a$ or $b$
approach zero while $N$ grows may reveal additional transition scales.
They cannot be inferred from the fixed-pair limit $R_N\to0$ alone.

The scalar proof suggests a broader approximation problem: characterize
families $f$ whose fixed-midpoint divided differences are monotone and
whose endpoint secants yield optimal multiplicative-increment bounds.
For Euler-type transformations, this could produce universal constants
for positive kernels beyond the harmonic double polylogarithms.

\subsection{Effective all-ratio moment asymptotics}

The theorem proves the existence of a remainder constant for each fixed
order. Produce explicit numerical constants valid for all $0\le m\le n$
above a stated $n_0$. This requires quantitative bounds on the normalized
phase and amplitude derivatives, plus the global phase-gap constant.
The existing interval proof of the slope inequality provides a concrete
starting point rather than an abstract compactness argument.

Can the inherited monotonicity of $J$ be proved by a short analytic
inequality, eliminating the compact interval computation? Ordinary
convexity of $\log\Gamma$ is insufficient, as the source already explains.
A successful proof would make the all-ratio theorem fully analytic and
could sharpen its uniform curvature bounds.

Find a uniform approximation retaining the exponentially small
fixed-$m$ corrections while also handling the balanced region. The
current theorem controls every algebraic order, but matching it to the
residue expansion requires analysis of the correction
$q(y)e^y/\gamma-1$ when $m e^{-s}$ changes scale. One target is an
explicit overlap theorem that recovers both expansions with their own
natural error budgets.

The probabilistic corollary invites quantitative questions. Establish a
uniform Kolmogorov or total-variation error bound for
$\sqrt H(Y-s)$, then a uniform Edgeworth expansion with tail control.
The metric, centering and normalization must be fixed before a rate is
claimed. Related questions concern logarithmic derivatives of $M_{n,m}$,
which encode the mean and covariance of $\log f(x)$ and $\log f(1-x)$
under the tilted moment measure.

\subsection{Appell mesh, saturation, and growing zero configurations}

The bound $\operatorname{mesh}(Q_n)\ge(3n)^{-(n-2)}$ is deliberately
coarse. Determine the actual asymptotic size of the smallest gap. Even a
polynomial lower bound $\operatorname{mesh}(Q_n)\gg n^{-A}$ would, through
the proved transfer argument, give a polynomial sufficient threshold
$k\gg n^{2A+2}$. Such an improvement would extend the growing-index regime
far beyond $n=O(\log k/\log\log k)$.

The current proof estimates the product over roots by its worst distance
to any root. Replace that bound with local reciprocal-gap sums such as
\[
 \sum_{\ell\ne j}|x_{n,j}-x_{n,\ell}|^{-1}
 \quad\hbox{and}\quad
 \sum_{\ell\ne j}|x_{n,j}-x_{n,\ell}|^{-2}.
\]
Together with centered gamma moments, this may improve the location
error $O(n/\sqrt{k})$ and recover constant offsets uniformly for growing
$n$. It would also distinguish extreme roots from interior roots when
their local geometries differ.

For each small index, determine the least $k$ that guarantees saturation
for every $\rho\in[0,1]$, and the least $k$ for each global motion law.
Certified continuation, exact endpoint analysis and the global
multiplicity bound should be used together. The new terminal-branch
theorem settles its direction whenever that branch is simple, but it
does not order the motion of all interior branches at small $k$.

More generally, replace $W_\rho$ by a positive family with a uniformly
bounded logarithmic derivative. The sign-transfer proof then changes
only through that bound. Determine which non-Lerch deformations inherit
effective saturation, and which require additional endpoint integrability.
This would isolate the root-stability mechanism from its original
special-function setting.

\subsection{Integral presentations and exact Gaussian identities}

The immediate identity target remains a finite rational certificate for
the existing $S_6$ formula. Use the integral shuffle diagonalization to
remove redundant coordinates before building Cayley, distribution,
conjugation and regularized double-shuffle rows. Each regularized row
must carry an analytic justification. A successful computation should
return an exact row combination producing the target, not only a rank
or a floating-point null vector.

\begin{conjecture}[A uniform Gaussian span question]
\label{conj:uniform-span}
For every integer $m\ge1$, the number
$S_{2m}+2\beta(2m)\log2$ belongs to the rational span of
\[
 \{g_{2m-2j,\,2j+1}:0\le j<m\}
 \cup\{\pi^{2m+1}\}
 \cup\{\beta(2j)\zeta(2m+1-2j):1\le j<m\}.
\]
\end{conjecture}

The proved $S_2,S_4$ cases in the repository and the conjectural
$S_6,S_8,S_{10},S_{12}$ relations motivate this formulation. Their
different evidential status is essential. The all-weight matrix theorem
justifies the chosen coordinates but does not prove that a mixed sum
belongs to their span. An appropriate generating-function identity could
settle many weights at once; until it is found, exact certificates at
successive weights are meaningful intermediate goals.

Determine analogous integral diagonal forms at even weights, other roots
of unity, and greater depth. The odd-weight proof uses the reflection
$y\leftrightarrow1-y$ and central factorials in a precise way; it does
not automatically furnish a higher-depth presentation. Integral data
would identify torsion primes and optimal denominator budgets before
expensive exact linear algebra is attempted.

\subsection{A practical order of attack}

The nearest finite targets are the $S_6$ row certificate, sharper
certified Appell mesh bounds, and explicit constants in the first-order
all-ratio moment error. The more structural targets are an analytic
global slope proof, a general mixed-Gaussian generating identity, and
an asymptotic theory for the Appell mesh. These tasks can proceed in
parallel because their proof dependencies are largely separate.

For every proposed identity, freeze the basis and integer vector before
validation; recompute it by an independent representation; produce exact
residual bounds; and pursue an algebraic or analytic certificate of
equality. For every asymptotic claim, state all joint parameter ranges
and preserve the distinction between finite diagnostics and a uniform
remainder proof. The present package supplies examples of both successful
proofs and a rigorously rejected numerical candidate.
